\section{Introduction}
\label{sec:intro}

Many scientific analyses share one bottleneck. Gold-standard labels are expensive, a much larger set of covariates is cheap, and a pretrained predictor can be run on every unit. Examples are remote-sensing land-cover classes, galaxy morphology, and census-based income or insurance indicators. The analyst wants a confidence interval for a population quantity such as a mean or prevalence, with guaranteed coverage and as little width as possible.

Prediction-powered inference (PPI) answers this with a rectifier: average the predictor over the unlabeled sample and add the labeled residual correction \citep{angelopoulos2023ppi}. PPI++ replaces the fixed correction with a data-driven weight \citep{angelopoulos2023ppiplus}, and cross-fitting handles predictors trained on the labels \citep{zrnic2023cross}. These methods all hold the data-collection design fixed in one respect: they accept the labeled set, the unlabeled set, and the predictor as given. Our question is what is left to gain under exactly that constraint. Holding labels, unlabeled covariates, and the predictor fixed, can the existing intervals be shortened without weakening coverage? If so, by how much, and under what conditions does the answer change?

We make three contributions.

\textbf{(1) An exact gap.} For the estimator class in which the mean is corrected by a calibrator $\hat m$ fitted on labeled data, the asymptotic variance is $V^\star+(1-a)\mathrm{Var}(\hat m-m)/n$, where $V^\star$ is the efficient bound for the information used, $m$ is the conditional mean, and $a=n/(n+N)$ is the labeled fraction (\cref{prop:var,thm:bound}). PPI++ equals the bound plus $(1-a)\Delta_L/n$, where $\Delta_L$ is the squared misfit of the best linear calibrator. PPI++ is therefore not efficient exactly when $m$ is nonlinear in the information used (\cref{cor:when}). The same argument shows that covariates beyond the predictor help through the same term (\cref{cor:cov}).

\textbf{(2) A practical condition.} A flexible calibrator helps only when its error variance is below the linear misfit, $\mathrm{Var}(\hat m-m)<\Delta_L$ (\cref{cor:learned}). This is a testable statement about the labeled sample, and it implies that flexibility can backfire at small $n$.

\textbf{(3) Evidence on known-truth and real data.} Using cross-fitting, the estimator stays exactly unbiased for any learner (\cref{prop:unbiased}). On synthetic problems, a spline calibrator shortens intervals by 17--19\% when calibration is strongly nonlinear, against an efficient-bound reduction of 24\%. It widens them when calibration is linear, by up to 16\% at $n=50$. On seven public pools with exact coverage checks, the estimated nonlinear share is at most 5\% of the linear residual variance, so there is little to gain. The flexible estimator is 5--11\% wider at $n=50$, 1--4\% wider at $n=200$, and within 1\% of PPI++ by $n=1000$. Coverage is at the nominal level for all methods, except for expected small-sample undercoverage on heavy-tailed incomes.

The answer is therefore conditional. The interval can be shortened with the same data and predictor only when the calibration function is materially nonlinear and the labeled sample is large enough to learn it. On the public pools we could check, the gain is not available at the sample sizes that matter most in practice. We state the conditions precisely, give the formulas to check them, and report the failures.
